\begin{tikzpicture}[font=\figurefont\small,
  box/.style={draw=BookRule,fill=BookPaper,align=center,
    text width=27mm,minimum height=14mm},
  arr/.style={-{Stealth[length=2mm]},draw=BookInk,line width=.8pt},
  sim/.style={arr,dashed,draw=BookTeal}]
  \node[box,fill=FigureInput!12] (env) at (0,0) {真实环境\\一步交互};
  \node[box,fill=FigureInput!12] (data) at (4.4,0) {真实记录\\$(s,a,r,s',c)$};
  \node[box,fill=FigureModel!12] (q) at (11,0) {统一Q更新\\$r+\gamma c\max Q$};
  \node[box,fill=FigureModel!12] (model) at (4.4,-2.3) {经验模型\\后继与奖励计数};
  \node[box,fill=FigureOutput!12] (sample) at (11,-2.3) {模拟记录\\$(\widetilde s,\widetilde a,\widetilde r,\widetilde s',\widetilde c)$};
  \draw[arr] (env) -- (data);
  \draw[arr] (data) -- (q) node[midway,above] {直接学习};
  \draw[arr] (data) -- (model) node[midway,left] {拟合};
  \draw[sim] (model) -- (sample) node[midway,above] {采样$n_{\rm plan}$次};
  \draw[sim] (sample) -- (q);
  \draw[arr] (q.north) -- (11,1.8) -- (0,1.8) -- (env.north);
  \node[fill=white] at (5.5,1.8) {按更新后的Q探索，选下一个真实动作};
  \node[align=center,text width=100mm] at (6,-4)
    {重用记录即可完成部分相同更新；\\查询新后果的能力来自模型，也承担模型误差。};
\end{tikzpicture}
